\documentclass[aspectratio=169,11pt]{beamer}

\usetheme{Madrid}
\usecolortheme{default}
\usefonttheme{professionalfonts}
\setbeamertemplate{navigation symbols}{}
\setbeamertemplate{footline}[frame number]

\usepackage{amsmath, amssymb, booktabs, array}

\title{Attention, Salience, Complexity, and Learning}
\subtitle{Behavioral Labor -- Week 4}
\author{Teaching deck draft}
\date{}

\begin{document}

\begin{frame}
  \titlepage
\end{frame}

\begin{frame}{Central question}
\begin{itemize}
    \item How do attention, salience, complexity, learning, and costly information acquisition shape labor choices?
    \item Week 3 studied beliefs and expectations in job search.
    \item Week 4 asks what happens when the relevant labor rule exists but is hard to notice, understand, or update.
    \item Core margins: earnings, hours, take-up, claiming, effort, saving, training.
\end{itemize}
\end{frame}

\begin{frame}{Bridge from Week 3}
\small
\begin{tabular}{p{0.30\linewidth} p{0.58\linewidth}}
\toprule
Week 3 object & Week 4 shift \\
\midrule
Subjective beliefs & Workers must notice and process relevant signals \\
Job search information & Broader labor schedules, contracts, benefits, and programs \\
Perceived returns & Perceived tax, benefit, bonus, saving, and training maps \\
Updating over unemployment & Learning from repeated payroll, policy, and workplace exposure \\
\bottomrule
\end{tabular}
\end{frame}

\begin{frame}{DellaVigna taxonomy: Week 4 highlighted}
\small
\begin{tabular}{p{0.31\linewidth} p{0.57\linewidth}}
\toprule
Branch & Labor interpretation \\
\midrule
Nonstandard preferences & Present bias, reference points, reciprocity, fairness \\
Nonstandard beliefs & Subjective expectations about wages, job finding, returns \\
\textbf{Nonstandard decision-making} & \textbf{Attention, salience, complexity, limited consideration, learning} \\
Responses & Firms and policies redesign the information environment \\
\bottomrule
\end{tabular}

\vspace{0.5em}
\textbf{This week:} workers respond to the schedule they perceive and learn, not only the schedule written down.
\end{frame}

\begin{frame}{Benchmark transparent labor choice}
\[
a_i^\star \in \arg\max_{a \in \mathcal{A}} U_i(a;\theta_i)
\quad \text{s.t.} \quad c_i = y_i + w_i a - T_i(a)
\]

\small
\begin{tabular}{p{0.17\linewidth} p{0.70\linewidth}}
\toprule
Object & Interpretation \\
\midrule
$a$ & hours, earnings, effort, claiming, saving, or training \\
$T_i(a)$ & true tax-benefit, payroll, or contract schedule \\
$U_i$ & payoff from consumption, work, effort, time, or program participation \\
Assumption & the worker understands the local action-to-payoff map \\
\bottomrule
\end{tabular}
\end{frame}

\begin{frame}{Salience wedge: perceived incentives}
\[
a_i \in \arg\max_{a \in \mathcal{A}} U_i(a;\theta_i)
\quad \text{s.t.} \quad \tilde{T}_i(a) \neq T(a)
\]

\begin{itemize}
    \item The worker responds to the perceived schedule $\tilde{T}_i(a)$.
    \item Low salience: the relevant incentive is not prominent at the moment of choice.
    \item Misperception: the worker holds the wrong local slope, threshold, or eligibility rule.
    \item Heuristic response: average incentives replace marginal incentives.
\end{itemize}
\end{frame}

\begin{frame}{Salience versus complexity versus information}
\scriptsize
\begin{tabular}{p{0.20\linewidth} p{0.25\linewidth} p{0.30\linewidth}}
\toprule
Object & Distortion & Labor signature \\
\midrule
Low salience & incentive exists but is not prominent & stronger response when highlighted \\
Complexity or opacity & action-to-payoff map is hard to decode & mistakes persist despite disclosure \\
Missing information & worker lacks a rule or fact & response to letters or tutorials \\
Endogenous information & worker chooses attention effort & response depends on stakes and costs \\
Dynamic learning & representation evolves with exposure & wedges shrink or change over time \\
\bottomrule
\end{tabular}

\vspace{0.6em}
Do not infer attention from attenuation alone. Name the margin and the behavioral object.
\end{frame}

\begin{frame}{Dynamic learning and costly attention}
\[
a_{it}, m_{it} \in \arg\max \tilde{U}_{it}(a;\theta_i,b_{it}) - C(m_{it})
\quad \text{with} \quad
b_{i,t+1}=B(b_{it},m_{it},x_{it+1})
\]

\begin{itemize}
    \item $m_{it}$ is costly information acquisition or attention effort.
    \item $b_{it}$ is the worker's representation of the schedule or contract.
    \item $x_{it+1}$ includes experience, letters, payroll feedback, peers, and caseworkers.
    \item Policy and firms can change current salience and the speed of learning.
\end{itemize}
\end{frame}

\begin{frame}{Labor supply under nonlinear schedules}
\begin{itemize}
    \item EITC, taxes, benefits, and earnings tests create nonlinear local incentives.
    \item Observed margins: earnings, hours, employment, bunching, and adjustment timing.
    \item Chetty-Friedman-Saez: local knowledge and EITC earnings responses.
    \item Kostol-Myhre: labor supply responses to learning the tax and benefit schedule.
    \item Main caution: muted bunching can reflect information, adjustment costs, constraints, or true low elasticity.
\end{itemize}
\end{frame}

\begin{frame}{Take-up, claiming, and policy navigation}
\small
\begin{tabular}{p{0.26\linewidth} p{0.25\linewidth} p{0.33\linewidth}}
\toprule
Intervention & Margin & Mechanism caution \\
\midrule
Information letter & application or take-up & awareness versus salience \\
Reminder & filing or claiming timing & attention versus pressure \\
Simplified notice & completed claim & complexity versus hassle cost \\
Trusted messenger & program participation & information versus legitimacy \\
\bottomrule
\end{tabular}

\vspace{0.6em}
Bhargava-Manoli style designs identify treatment effects on claiming, but not a single mechanism without additional diagnostics.
\end{frame}

\begin{frame}{Opaque incentives and workplace effort}
\begin{itemize}
    \item Firms write piece rates, bonuses, thresholds, penalties, and performance metrics.
    \item Workers may know there is an incentive while misunderstanding the formula.
    \item Observed margins: effort, output, task allocation, bonus attainment.
    \item Abeler-Huffman-Raymond style designs vary complexity and observe effort response.
    \item Workplace opacity can be accidental, productivity-enhancing, or strategic.
\end{itemize}
\end{frame}

\begin{frame}{Savings, retirement, and training}
\begin{itemize}
    \item Payroll-linked saving depends on defaults, disclosures, and contribution understanding.
    \item Social Security and retirement claiming require mapping current work to future benefits.
    \item Training decisions require learning returns, subsidies, deadlines, and procedural steps.
    \item Observed margins: participation, contributions, claiming, enrollment, completion.
    \item Same diagnostic habit: what rule is perceived, what action moves, what remains confounded?
\end{itemize}
\end{frame}

\begin{frame}{Methods and identification map}
\scriptsize
\begin{tabular}{p{0.28\linewidth} p{0.28\linewidth} p{0.30\linewidth}}
\toprule
Design & Best identified object & What it cannot separate alone \\
\midrule
Information letter or tutorial & missing information response & salience, trust, reminders \\
Simplification or disclosure & complexity and legibility & hassle costs, salience \\
Schedule change plus exposure & learning and attenuated response & sorting and diffusion \\
Monitored attention & costly information acquisition & experiment behavior \\
Workplace incentive experiment & effort under opacity & morale, peer effects, reciprocity \\
Structural learning model & welfare and counterfactuals & functional-form dependence \\
\bottomrule
\end{tabular}
\end{frame}

\begin{frame}{Welfare and design}
\begin{itemize}
    \item Stable inattention: observed choices may poorly reveal welfare.
    \item State-dependent inattention: workers may err after new rules, transitions, or shocks.
    \item Learning: short-run wedges and long-run wedges can differ sharply.
    \item Endogenous attention: optimal policy may reduce information costs rather than change payoffs.
    \item Design question: who makes the environment legible, and who benefits from opacity?
\end{itemize}
\end{frame}

\begin{frame}{Week 4 lab}
\begin{itemize}
    \item Reproduce: synthetic Kostol-Myhre-style schedule-learning factbook.
    \item Diagnose: true incentives, perceived incentives, exposure, information, and hours response.
    \item Transfer: Bhargava-Manoli-style policy navigation with reminders and simplification.
    \item Optional extension: Abeler-Huffman-Raymond-style workplace incentive opacity.
\end{itemize}
\end{frame}

\begin{frame}{Bridge to Week 5}
\begin{itemize}
    \item Week 4: workers respond to perceived and learned schedules.
    \item Week 5: firms design incentives, contracts, monitoring, and workplace communication.
    \item The bridge is contract legibility.
    \item Once workers are bounded, employer design becomes part of the labor mechanism.
\end{itemize}
\end{frame}

\begin{frame}{Takeaways}
\begin{itemize}
    \item Attention, salience, complexity, and learning are dynamic labor mechanisms.
    \item The same reduced-form response can mix information, salience, hassle, optimization frictions, and learning.
    \item Labor applications must name the observed margin before naming the behavioral object.
    \item Welfare depends on persistence, heterogeneity, endogenous attention, and institutional design.
\end{itemize}
\end{frame}

\end{document}
